\begin{document}
\twocolumn[
  \begin{@twocolumnfalse}
  \maketitle
  \begin{abstract}
  \noindent
  The Israel-Hamas conflict has unfolded not only through military and diplomatic
  events but also through continuous public argument online. We analyse
  \textbf{1{,}382{,}896} stance-labelled comments - \textbf{1{,}004{,}629} from Reddit
  and \textbf{378{,}267} from YouTube - over a \emph{matched}
  7~October~2023 to 7~May~2024 window, to ask how platform architecture shapes
  political discourse. Three findings are not predictable from platform design
  alone. First, emotional tone is \emph{politicised} on Reddit, where both partisan
  camps are far more negative than neutral comments, but only \emph{weakly coupled}
  to stance on YouTube; the contrast holds under two sentiment instruments and on
  human labels. Second, negativity
  and toxicity are related but not interchangeable: most toxic Reddit comments
  are negative, yet only about one negative comment in four is toxic, and the two
  partisan camps are equally toxic while both exceed neutral comments.
  Third, Reddit's \emph{reply} network is \emph{disassortative} by stance
  (Newman $r=-0.184$ against a degree-preserving null centred on zero): 62.4\% of
  partisan replies cross the divide. This replicates, in a geopolitical-conflict
  corpus, the cross-cutting interaction already reported for US party politics; what
  is new is that the \emph{same users} appear strongly clustered when measured by
  thread co-membership, which locates the long-running echo-chamber disagreement in
  the choice of measure rather than in the platform. Stance classifiers transfer poorly across
  platforms, and we show this survives length matching, so it reflects vocabulary
  rather than comment length. We release a reproducible LLM-assisted labelling
  pipeline together with a three-annotator human validation (Krippendorff's $\alpha=0.796$), which finds stance accuracy of $0.836$ on Reddit but $0.618$ on YouTube. We subject the platform contrast to three adversarial checks built from that validation - degrading Reddit's labels to YouTube's measured quality, removing off-topic comments with a validated relevance classifier, and correcting both platforms for their measured misclassification - and it survives all three.
  \vspace{0.6em}
  \end{abstract}
  \end{@twocolumnfalse}
]

% ============================================================
\section{Introduction}
% ============================================================
%\subsection{Research Motivation}
The Hamas-led attack on southern Israel of 7~October~2023 and the Israeli military
campaign in Gaza that followed became one of the most heavily discussed
geopolitical events of the social-media era. Ng and Chow~\cite{ngchow2025} index
roughly 1.08~billion English-language posts across Twitter/X, Reddit and TikTok in
the first four months alone, 59\% of them negative in sentiment, and identify eight
distinct volume peaks driven both by battlefield events and by offline protest and
boycott movements. Public argument about the conflict was therefore not a marginal
accompaniment to it; for most audiences outside the region it was the primary
channel through which the conflict was encountered.

That argument also changed the character of the forums carrying it. Hofmann
et al.~\cite{hofmann2026} annotated YouTube comments on the conflict and found
hate speech in \textbf{40.4\%} of those beneath public-broadcaster videos and
\textbf{31.6\%} beneath private ones, against a stance distribution that was
predominantly neutral in both. Rosen and Walther~\cite{rosen2025} document a
parallel escalation in anti-Jewish and anti-Muslim posting in the month after
7~October, and show that the replies such posts attract feed back into how
hatefully their authors post next.
Comment threads that ordinarily host reaction became a venue for
sustained partisan dispute, which is precisely the material this paper analyses:
not what participants believe, but how the argument is organised once a platform's
design decides who can reply to whom, at what length, and with what visibility.

Because platforms differ in architecture, they may shape that argument in
systematically different ways. We compare two with sharply contrasting designs:
\textbf{Reddit}, a threaded forum with voting and unlimited comment length, and
\textbf{YouTube}, a media-centric platform where flat, length-limited comments react
to video content.

We treat the structural contrast between these designs as a \emph{premise and a
control}, not as a result. That Reddit hosts longer threaded exchanges while YouTube
hosts shorter reactions is a well-documented property of the two platforms and is
not a finding about this conflict. Our question is what remains once that difference
is held in view: whether the \emph{relationships} among stance, emotion, toxicity and
interaction differ across platforms in ways architecture alone does not predict.

\subsection{Research Questions}
Because sentiment and toxicity are likely to overlap, we study them jointly rather
than as separate questions, giving three:

\begin{itemize}[leftmargin=1.2em,itemsep=1pt,topsep=2pt]
  \item \textbf{RQ1 (Affective tone and toxicity).} How do emotional tone and
  toxicity differ between platforms and stances, and are they the same phenomenon?
  \item \textbf{RQ2 (Topics and vocabulary).} What distinct topics and lexical
  choices characterise each platform and stance?
  \item \textbf{RQ3 (Polarization and interaction structure).} Do users cluster into
  same-stance interaction, and is stance encoded in platform-specific language?
\end{itemize}

\subsection{Contributions}
This work contributes (1)~large-scale evidence that the \emph{coupling} between
stance and emotional tone is platform-dependent - strong on Reddit, weak on
YouTube; (2)~a quantification of how far toxicity and negativity
diverge in this corpus - the distinction is recognised in principle, but we measure
the overlap directly across the full corpus and show that the two partisan camps
are equally toxic; (3)~evidence that thread co-membership and actual reply
behaviour give \emph{opposite} answers on the same users, which reframes the
long-running echo-chamber disagreement as substantially a question of measurement; and (4)~a reproducible LLM-assisted labelling
pipeline with retained per-label justifications, a completed three-annotator human validation of it and of the sentiment instruments, and three reusable procedures for bounding a cross-platform finding by the annotation quality behind it, together with the released validation
protocol.

Cross-platform stance-transfer failure is \emph{not} claimed as novel - it is
established for stance datasets generally by Ng and Carley~\cite{ng2022}. Our
contribution there is narrower: the gap persists after controlling comment length,
and we identify which lexical families fail to carry across.

% ============================================================
\section{Related Work}
% ============================================================
\subsection{Platform Architecture and Political Discourse}
Social media platforms differ in structural affordances that systematically
shape user behaviour. Boulianne et al.~\cite{boulianne2025} compare seven
platforms, including Reddit and YouTube, using survey data from four countries,
and distinguish them by their \emph{networking features}: structural ones such
as user connectivity, privacy and anonymity, and social ones such as perceived
audience and the norms of interaction. They report that political ideology
predicts both adoption and political use across all seven. Shugars and Ha~\cite{shugars2025}
compared Twitter and Reddit directly, finding that Reddit's \emph{conversational
visibility} - the ability for users to see and engage with exchanges as a
whole - enabled more cross-cutting and productive political discourse than
Twitter's broadcast-oriented design. The present study extends this
platform-comparative framing to YouTube, which shares neither Reddit's deep
threading nor Twitter's follow-graph, and applies it to a substantially larger corpus than
prior comparisons have used.



\subsection{Online Discourse on the Israel-Hamas Conflict}
The October 2023 escalation generated a rapid surge of multilingual digital
discourse that researchers have begun to quantify computationally. Antonakaki
et al.~\cite{antonakaki2025} analysed Telegram, Twitter/X, and Reddit from October
2023 to mid-2025 using LDA, BERTopic, and transformer-based emotion models,
identifying platform-specific narrative strategies and propaganda amplification
patterns. Guerra et al.~\cite{guerra2025} focused on Reddit, scoring over
450{,}000 posts from four conflict-related subreddits on a lexicon-based index of
anger, polarity and subjectivity, and found peaks in that index aligning with
specific events on the ground. Closest
to our own design, Hofmann et al.~\cite{hofmann2026} hand-annotated 4{,}983
YouTube comments on the conflict for hate speech and for the same three stances
we use, trained classifiers reaching AUROC 0.83 to 0.90, and applied them to
comment sections from public and private news channels. Their corpus is roughly
two orders of magnitude smaller than ours and they do not compare platforms, but
their stance scheme and their finding of a predominantly neutral distribution
both anticipate ours. Notably, Antonakaki et al.\ already report that Reddit
hosts the more reflective and deliberative discussion of the platforms they
compare, which is one reason we treat that contrast as an established premise
rather than a finding of ours.

These studies establish that the conflict produced intense and polarized digital
discourse, but they leave three gaps this work addresses. First, the
cross-platform comparisons of this conflict omit \textbf{YouTube}, so the
media-centric reaction layer is absent from them. Second, none assigns
\textbf{per-comment stance} at corpus scale, which is what allows sentiment, toxicity
and interaction to be analysed \emph{conditional on} political position. Third, none
examines \textbf{interaction structure} - reply networks, homophily, or
cross-platform stance transfer - so the question of whether these communities
actually talk to each other is untouched.



\subsection{Sentiment Analysis of Political Text}
Lexicon-based sentiment methods remain standard for large-scale political corpora.
VADER~\cite{hutto2014}, designed for microblog contexts, captures punctuation
intensity, slang, and negation. TextBlob's default analyser is also lexicon-based:
it matches tokens against a hand-curated dictionary assigning each entry a polarity
in $[-1,1]$, a subjectivity in $[0,1]$ and an intensity weight, then averages the
matched polarities across a sentence while applying negation and intensifier
modifiers. Polarity is therefore a signed average of word-level priors, while
subjectivity measures how opinion-laden (not how positive) the wording is; we
classify a comment positive above $+0.1$ and negative below $-0.1$. Because both
tools are lexicon-based and neither is trained on this domain, we additionally apply
a RoBERTa model fine-tuned on Twitter data~\cite{loureiro2022} - building on the
BERT architecture of
Devlin et al.~\cite{devlin2019} - and treat agreement among the three as a
measurement question in its own right rather than an assumption. Because agreement
among instruments does not establish which of them is right, we also validate all
three against human labels (Section~\ref{sec:sentval}).

\subsection{Topic Modeling}
Latent Dirichlet Allocation~\cite{blei2003} and Non-negative Matrix
Factorization~\cite{leeseung2000} represent documents as mixtures over latent
topics; LDA assumes a Bayesian generative process while NMF performs a matrix
decomposition that tends to produce sparser topics. BERTopic~\cite{grootendorst2022}
instead clusters sentence-embedding representations. We apply LDA and NMF in
parallel so that no single algorithm drives the reported structure, and quantify
topic quality with $c_v$ coherence. We emphasise
that topic models recover \emph{co-occurring vocabulary}, not narratives: in conflict
research a narrative is a causally linked account assigning agency, responsibility
and moral valence, and recovering that requires statement-level analysis we do not
attempt here.

\subsection{Echo Chambers and Partisan Polarization}
The homophily principle - that similar people interact at higher rates than
dissimilar ones~\cite{newman2003} - underpins computational echo-chamber
research. Where filtering and self-selection are expected to narrow exposure,
Bail et al.~\cite{bail2018} showed experimentally that
exposure to opposing views can \emph{increase} polarization. Cinelli
et al.~\cite{cinelli2021} demonstrated that echo-chamber strength is itself
platform-dependent, motivating our comparative design. Most directly relevant to
our RQ3, De Francisci Morales et al.~\cite{defrancisci2021} reconstructed the
political interaction network of Trump and Clinton supporters on Reddit and found a
\emph{preference} for cross-cutting replies over within-group ones - asymmetrically
so - against a degree-preserving rewiring null. We therefore do not claim
cross-cutting interaction on Reddit as a new result. We ask instead whether it holds
in a geopolitical-conflict setting rather than domestic party politics, and why it
coexists with strong echo-chamber readings of the same platform. Crucially, homophily is a
property of a \emph{network}, so we measure it on an explicit user-user reply graph
and report Newman's assortativity coefficient for categorical
attributes~\cite{newman2003} alongside a chance-corrected homophily index; raw
same-stance rates are uninterpretable without the baseline expected under random
mixing.

\subsection{Toxicity}
The Perspective API operationalizes toxicity through the framework of Wulczyn,
Thain and Dixon~\cite{wulczyn2017}, scoring the probability that a rater would find a
comment rude or disrespectful across six attributes. We use classifiers of this
kind in preference to zero-shot toxicity prompting of a general-purpose LLM for three
reasons: their scores are calibrated against human rater distributions, they are
directly comparable to a large existing literature, and they are stable across runs,
whereas LLM zero-shot toxicity is prompt-sensitive, uncalibrated and not comparable
across studies. Their known demographic biases mean we report effect sizes rather than
thresholded counts wherever possible. Perspective is, however, a hosted service that
is being retired, so its scores cannot be reproduced. Our primary instrument is
therefore Detoxify's \texttt{unbiased} model~\cite{detoxify2020}, an openly
released RoBERTa classifier trained on the Jigsaw toxicity challenge data,
including the \emph{Unintended Bias in Toxicity Classification} benchmark built to
reduce identity-related bias, which can be run on the full corpus and re-run by
anyone; Perspective serves as a cross-check
on a subsample. Throughout we say ``toxicity'' and name the specific attribute; we
avoid the vaguer term ``harm'', which we never define independently of these
scores.

\subsection{LLM-Assisted Annotation}
Manually labeling millions of comments is impractical at this scale. Gilardi,
Alizadeh and Kubli~\cite{gilardi2023} showed that LLMs can outperform crowd workers
on stance and frame detection. We use Llama-3.3-70B-Instruct, retaining per-label
confidence tiers and free-text justifications. Because automated annotation without
human verification is not sufficient evidence of label quality, we specify a
protocol (Section~\ref{sec:validation}) reporting the annotator's test-retest reliability, the empirical value of the confidence filter, and a three-annotator human study measuring
model-human agreement.
Classification uses scikit-learn~\cite{pedregosa2011}; agreement statistics follow
Krippendorff~\cite{hayes2007}.

% ============================================================
\section{Data and Methodology}
\label{sec:data}
% ============================================================
\subsection{Data Collection}
Reddit comments were collected with the \textit{Arctic Shift} historical-dump tool
from five conflict-relevant subreddits - \texttt{r/AskMiddleEast},
\texttt{r/IsraelPalestine}, \texttt{r/Israel}, \texttt{r/IsraelPalestineWar\_23},
and \texttt{r/Palestine} - yielding about \textbf{3.08 million} raw comments, each
joined to its parent post so thread context is available.

YouTube candidate videos were discovered through the \textit{YouTube Data API v3}
using roughly two dozen conflict-related queries restricted to the
7~October~2023 to 7~May~2024 window, returning about 2{,}600 unique videos. Comments
were then scraped from YouTube's HTML and AJAX continuation endpoints, producing
\textbf{789{,}328} raw comments.

\subsection{Temporal Anchoring and the Matched Window}
Two timestamp issues govern every temporal comparison. Reddit comments carry true
posting timestamps. YouTube comment times, however, were reconstructed from relative
``x ago'' strings at scrape time; because those strings are coarse at the top of
their range (any comment 24 to 35 months old renders as ``2 years ago'') and no exact
per-comment scrape time was persisted, the reconstructed field spans 2024 to 2026 and
is not a usable comment timeline. We therefore anchor YouTube on the \emph{video
publish date}, and note that recovering exact \texttt{publishedAt} values through the
Data API is the appropriate long-term fix.

Because sentiment, toxicity and polarization can all drift over time, comparing
platforms across different periods would confound platform with period. We therefore
restrict \emph{both} platforms to the common window
\textbf{7~October~2023 to 7~May~2024}, removing a Reddit-only pre-war baseline of
\textbf{42{,}249} comments (3.81\%) and leaving YouTube unchanged. Figure~\ref{fig:temporal} shows
the resulting monthly volume on both platforms over that window.

\begin{figure}[H]
  \centering
  \includegraphics[width=\linewidth]{f11_temporal.png}
  \caption{Monthly comment volume over the matched window. Reddit is anchored on the comment timestamp, YouTube on the video publish date.}
  \label{fig:temporal}
\end{figure}



\subsection{Stance Labeling}
Both platforms were annotated by the \emph{same} automated annotator, so the stance
scheme is consistent across them. We used \texttt{Llama-3.3-70B-Instruct} with an
expert-annotator system prompt embedding the full guidelines. Each item was presented
with its context, and the model returned a stance label, a confidence level
(High/Medium/Low) and a justification retained for auditability, on a four-class
scheme: Pro-Palestine~(P), Pro-Israel~(I), Neutral~(N) and Irrelevant~(R). Figures abbreviate the two partisan classes as \emph{Pro-Pal} (Pro-Palestine) and \emph{Pro-Isr} (Pro-Israel). Labelling
ran asynchronously in 50-comment batches at temperature $0.1$ with per-batch
checkpointing.

\subsection{Label Validation}
\label{sec:validation}
Label quality has two separable components. \emph{Reliability} asks whether the
annotator gives the same answer twice; \emph{accuracy} asks whether that answer is
right. We measure both: reliability on the labelled corpus itself, and accuracy
against three human annotators working blind.

\subsubsection*{Test-retest reliability}
Comments were submitted in 50-item batches in file order, so two occurrences of an
identical comment text more than 50 rows apart were annotated in separate batches
and are independent re-annotations of the same item. Restricting further to repeats
within the \emph{same} post or video holds the surrounding context fixed. On the
High-confidence labels that enter the analytic corpus the annotator reproduces its
own label on \textbf{70.4\%} of 1{,}981 repeated Reddit items ($\kappa = 0.537$) and
\textbf{60.9\%} of 1{,}295 repeated YouTube items ($\kappa = 0.410$). Pooling all
confidence tiers lowers this to $\kappa = 0.522$ on Reddit and $0.254$ on YouTube,
so the stated confidence does track label stability.

\subsubsection*{Human validation}
Three annotators independently labelled a stratified sample of \textbf{270}
comments, balanced across platform $\times$ stance $\times$ confidence tier (30 High
and 15 Medium/Low per stance per platform) drawn from the \emph{pre-filter} labelled
data. Annotators saw the comment and its context but never the model's label, and
worked from the same written guidelines given to the model. The gold standard is the
majority vote.

The annotators agree substantially with each other: Krippendorff's
$\alpha = \textbf{0.796}$, Fleiss' $\kappa = 0.796$, pairwise Cohen's $\kappa$
between $0.761$ and $0.836$, and all three choose the same label on \textbf{78.1\%}
of items. Agreement is comparable on both platforms (mean pairwise $\kappa$ $0.813$
Reddit, $0.774$ YouTube). The task is therefore well posed and the gold standard is
not itself the weak link.

Measured against that standard the model is weaker than its confidence tiers
suggest (Table~\ref{tab:human}). Two results matter and neither is comfortable.

\textbf{The confidence filter does not screen out irrelevance.} Across the sample,
\textbf{17.0\%} of comments the model assigned a stance are judged \emph{irrelevant}
to the conflict by the human majority, and that rate is essentially identical in the
retained High tier (17.2\%) and the discarded Medium/Low tier (16.7\%). Roughly one
in six comments in the analytic corpus is therefore off-topic content carrying a
spurious stance label. The filter does work for its other purpose: among items whose
gold label is a stance, High-tier accuracy is $0.600$ against $0.433$ for
Medium/Low, a gap of $+16.7$ points that justifies keeping the filter.

\textbf{Label quality is markedly worse on YouTube.} On High-tier items that humans
judge to carry a stance at all, the model matches the human gold standard on
\textbf{83.6\%} of Reddit comments but only \textbf{61.8\%} of YouTube ones. The
YouTube errors are concentrated in the Neutral class, which the model frequently
reads as Pro-Israel.

This gap has an identifiable cause in our own pipeline rather than in the
platforms. The Reddit prompt supplied the subreddit, the parent post's title and
its body alongside each comment; the YouTube prompt supplied only an opaque video
identifier, the author handle and a like count. YouTube comments were therefore
annotated without the discursive context that Reddit comments had, and the
instruction to weigh ``video context'' was issued without that context being
present, which plausibly encouraged the model to supply one: a recurring pattern
in the YouTube errors is a justification describing material absent from the
comment. We therefore attribute much of the $0.836$ versus $0.618$ gap to this
asymmetry rather than to an intrinsic difficulty of YouTube text, and note that
our human annotators worked under the same handicap, seeing only a video
identifier themselves, and still agreed among themselves at $\alpha = 0.796$. Short,
multilingual and formulaic comments remain the harder annotation problem, but they
are not the whole of it.

\subsubsection*{What video context is worth: a controlled test}
We test that attribution rather than assert it. Taking the 131 YouTube validation
items on which a human majority exists, we re-annotate them twice under identical
conditions, using the same model, provider, guidelines, temperature and batch size,
and varying only whether each item carries its video title and description. Both
arms use a different open-weight model from the deployed annotator, so what the
comparison isolates is the effect of context rather than of model choice.

The no-context arm reproduces the corpus result exactly: $0.618$ against the human
gold standard on High-confidence stance items, the same figure the deployed
pipeline achieves. The shortfall is therefore a property of the prompt rather than
of the annotator model.

Adding video context raises stance accuracy from $0.651$ to $\textbf{0.734}$ on
items humans judged to carry a stance, improving every class (Pro-Palestine $+12.5$
points, Pro-Israel $+8.0$, Neutral $+5.8$). It simultaneously degrades relevance
judgement: the share of human-irrelevant comments wrongly assigned a stance rises
from \textbf{18\%} to \textbf{73\%} (Table~\ref{tab:context}). Told that the video
concerns the conflict, the model infers that the comment must concern it too.

\begin{table}[t]
  \centering
  \small
  \caption{What video context is worth. Both arms use the same model, prompt and
  settings on the same 131 items; only the video title and description differ.}
  \label{tab:context}
  \setlength{\tabcolsep}{4pt}
  \begin{tabular}{@{}lcc@{}}
    \toprule
    & Stance accuracy & Irrelevant given \\
    & ($n=109$) & a stance ($n=22$) \\
    \midrule
    No video context   & 0.651 & 18\% \\
    With video context & 0.734 & 73\% \\
    \bottomrule
  \end{tabular}
\end{table}

This carries a design implication, and it is why we recommend separating the two
decisions rather than merely supplying more context. Relevance and stance want
video context in opposite directions: relevance is better judged on the comment
alone, stance with the video in view. A single four-way prompt cannot optimise
both, which also explains the retained irrelevance reported above. These figures
rest on 131 items, 22 of them irrelevant, and a single run; we report them as a
diagnosis of our own pipeline rather than as a general benchmark.

\begin{table}[t]
  \centering
  \small
  \caption{Human validation against a three-annotator majority gold standard
  ($n=270$, blind). Stance accuracy and retained irrelevance are reported on the
  High-confidence tier, which is the corpus actually analysed.}
  \label{tab:human}
  \begin{tabular}{lrr}
    \toprule
     & Reddit & YouTube \\
    \midrule
    Inter-annotator $\kappa$ (mean pairwise) & 0.813  & 0.774  \\
    Model accuracy, stance classes           & 0.750  & 0.562  \\
    \midrule
    \multicolumn{3}{l}{\emph{High-confidence tier}} \\
    Irrelevant content retained              & 18.9\% & 15.6\% \\
    Stance accuracy on the remainder         & 0.836  & 0.618  \\
    \bottomrule
  \end{tabular}
  \vspace{2pt}
  \footnotesize Pooled across platforms: Krippendorff's $\alpha = 0.796$;
  78.1\% of items unanimous; stance-class macro-F1 $0.656$.
\end{table}

\subsubsection*{Three adversarial checks on the headline result}
Both validation findings threaten the same claim: that stance and sentiment are
coupled on Reddit (Cram\'er's $V = 0.210$ with RoBERTa) and only weakly on YouTube
($V = 0.073$). Label
noise attenuates association, and Reddit's labels are the better ones, so the
contrast could restate the difference in label quality; off-topic comments add
noise of their own. We test these objections directly rather than discussing them
(Table~\ref{tab:robust}).

\textbf{Degrading Reddit to YouTube's label quality.} Using the error structure
measured on the YouTube half of the validation sample, we corrupt 26\% of Reddit's
stance labels so that Reddit's accuracy falls to YouTube's $0.618$, and recompute
the association over 200 simulations. Reddit's $V$ falls to $0.118$
$[0.117, 0.119]$ and remains \textbf{1.6 times} YouTube's, far outside the
simulated interval (with VADER, $0.082$ and $2.1$ times).

\textbf{Removing off-topic comments outright.} We train a relevance classifier on
the model's own Irrelevant labels, which are plentiful, and tune its threshold on
the human annotations, which are independent of the model. Against human
judgements it reaches ROC-AUC $0.831$ on Reddit and $0.744$ on YouTube. Applying
it removes 12.7\% of the Reddit corpus and 20.0\% of YouTube's. Reddit's $V$ falls
to $0.180$; YouTube's barely moves ($0.073$ to $0.072$), leaving the
contrast at \textbf{2.5 times}. The stance-conditional pattern is unchanged on both
platforms: on Reddit both partisan camps remain far more negative than Neutral
(Pro-Palestine $-0.635$, Pro-Israel $-0.606$, Neutral $-0.317$), while on YouTube
the three stances remain much closer together ($-0.258$, $-0.148$, $-0.136$). Recomputed identically before and
after, the RQ3 reply graph becomes \emph{more} disassortative once off-topic
comments are removed, and the share of cross-partisan replies rises from 61.6\% to
65.2\%, so that finding is not a product of the contamination either.

\textbf{Correcting both platforms for their measured error.} The complementary
analysis disattenuates rather than degrades. Treating each platform's measured
confusion matrix as a misclassification process and solving for the underlying
table raises Reddit to $V=0.293$ $[0.235, 0.430]$ and YouTube to $0.105$
$[0.056, 0.274]$, a ratio of $2.8$ (with VADER, $0.207$ and $0.054$). Both platforms
rise, the contrast persists, and the corrected YouTube value is small rather than
nil, which is how we describe it.
These intervals are wide because the confusion matrices rest on 73 and 76 items
respectively, and at their extremes they overlap; we therefore read this as
corroboration of the contrast's direction rather than a precise estimate of its
size.

\begin{table}[t]
  \centering
  \small
  \caption{Stance-sentiment association (Cram\'er's $V$) under three adversarial
  checks, with RoBERTa tone labels; the last column gives the Reddit/YouTube ratio
  with VADER. Neither differential label quality nor off-topic content accounts for
  the Reddit/YouTube contrast.}
  \label{tab:robust}
  \setlength{\tabcolsep}{3pt}
  \begin{tabular}{@{}lcccc@{}}
    \toprule
    & Reddit & YouTube & ratio & VADER \\
    \midrule
    As reported        & 0.210 & 0.073 & 2.9$\times$ & 3.8$\times$ \\
    Labels degraded    & 0.118 & 0.073 & 1.6$\times$ & 2.1$\times$ \\
    Off-topic removed  & 0.180 & 0.072 & 2.5$\times$ & 3.0$\times$ \\
    Error-corrected    & 0.293 & 0.105 & 2.8$\times$ & 3.8$\times$ \\
    \bottomrule
  \end{tabular}
\end{table}

We therefore report the magnitude of the contrast as an upper bound, and its
existence and direction as robust. We do not adopt the filtered corpus as the
primary one, for a reason worth stating: because \emph{neutral} and
\emph{irrelevant} are lexically similar, the filter removes Neutral comments
disproportionately (51\% of its Reddit removals, against Neutral's 18.6\% share)
and discards 11\% of Reddit and 25\% of YouTube comments that humans judge to carry
a genuine stance. That trades a known bias for a less legible one. The filter is a
robustness instrument here, and we release it as such.

\subsubsection*{What this means for the results}
We report the validation in full rather than only where it flatters the pipeline.
Three consequences follow. First, between 10\% (all three annotators agreeing) and
16\% (majority vote) of the analytic corpus is off-topic; this adds noise to every
aggregate, but at a similar rate on both platforms and, as shown above, without
changing any reported comparison. Second, stance-conditional effects on YouTube are
the least reliable quantities in this paper and we treat them as directional.
Third, the Reddit results, including the reply-network analysis carrying RQ3, rest
on labels that are $83.6\%$ accurate against a gold standard with $\alpha = 0.796$,
which is adequate for the population-level comparisons made here. Sampling code,
blind sheets, the three completed annotation sets, the relevance classifier and the
scoring scripts are released, as is an index of the 2{,}637 videos used in the context experiment above.

\subsubsection*{Validating the sentiment instruments}
\label{sec:sentval}
The same 270 comments were labelled for tone (positive, negative or neutral) by two
of the annotators working blind, each seeing the comment, its context and, for
replies, the comment being answered; the third annotator resolved the 36 items
(13.3\%) on which they disagreed. The two agree at Cohen's $\kappa = 0.78$ (Reddit
$0.79$, YouTube $0.76$). Against this standard \texttt{twitter-roberta} is the most
accurate instrument (accuracy $0.68$, $\kappa = 0.46$), ahead of VADER ($0.52$,
$\kappa = 0.27$) and TextBlob ($0.32$, $\kappa = 0.05$). We therefore report RoBERTa
as the primary tone measure and VADER alongside it.

The human labels reproduce the central RQ1 contrast: with human labels for both
stance and tone, the stance-tone association on High-confidence items is about twice
as strong on Reddit as on YouTube (Cram\'er's $V = 0.355$ against $0.175$). They do
not reproduce a platform difference in overall negativity. Reweighted to each
platform's stance composition, humans judge 56\% of Reddit and 61\% of YouTube
comments negative, whereas both RoBERTa and VADER score Reddit 15 to 21 points more
negative. The instruments appear to under-read negativity in YouTube's shorter,
more informal comments, so we treat the platform-level tone gap as
instrument-dependent and rest no claim on it. An earlier labelling round in which
annotators did not see the comment being replied to gave the same conclusions
($\kappa = 0.85$; RoBERTa again closest to the human labels; equal negativity on
both platforms; association 1.6 times stronger on Reddit).

\subsection{Corpus Cleaning and Final Datasets}
Two filtering decisions shape the analytic corpus. First, Irrelevant~(R) comments
and non-high-confidence labels are discarded. This removes only what the \emph{model}
called irrelevant; validation shows it misses roughly one comment in six, which is a
limitation of the filter rather than of the sample, and one an explicit relevance
classifier would address. Second - and unlike our earlier
analysis - automated accounts (\texttt{AutoModerator}, \texttt{*-ModTeam}) and the
collapsed \texttt{[deleted]} placeholder are removed \emph{corpus-wide} rather than
only from user-level analyses. These supply 5.93\% of the windowed Reddit corpus, and
their contamination is strongly asymmetric: 14.2\% of Neutral and 6.5\% of
Pro-Palestine comments but only 0.3\% of Pro-Israel comments, because subreddit rule
text is itself stance-inflected. Leaving them in inflated two stance classes and
produced a spurious ``community governance'' topic. YouTube shows no comparable
automated-account concentration.

The resulting corpus is \textbf{1{,}004{,}629} Reddit and \textbf{378{,}267} YouTube
comments (Table~\ref{tab:data}, Figure~\ref{fig:stance}). Stance shares carry
bootstrapped 95\% confidence intervals (4{,}000 resamples); at this scale the
intervals span roughly $\pm0.1$ percentage points, so sampling error is negligible
and the operative uncertainty is annotation error, which
Section~\ref{sec:validation} addresses.

\begin{table*}[t]
  \centering
  \small
  \caption{Corpus construction and final stance composition (matched window,
  automated accounts removed). Stance shares are given as the point estimate $\pm$ the half-width of a bootstrapped 95\% confidence interval (4{,}000 resamples).}
  \label{tab:data}
  \begin{tabular}{lrr}
    \toprule
    Stage & Reddit & YouTube \\
    \midrule
    Raw comments collected            & $\sim$3{,}083{,}006 & 789{,}328 \\
    Non-R, high-confidence            & 1{,}110{,}153 & 378{,}267 \\
    After matched window              & 1{,}067{,}904 & 378{,}267 \\
    After removing automated accounts & \textbf{1{,}004{,}629} & \textbf{378{,}267} \\
    \midrule
    Pro-Palestine (P) & 455{,}709 \quad 45.36\% $\pm$0.10\% & 212{,}217 \quad 56.10\% $\pm$0.16\% \\
    Pro-Israel (I)    & 362{,}340 \quad 36.07\% $\pm$0.10\% & 131{,}357 \quad 34.73\% $\pm$0.16\% \\
    Neutral (N)       & 186{,}580 \quad 18.57\% $\pm$0.08\% & \phantom{0}34{,}693 \quad \phantom{0}9.17\% $\pm$0.10\% \\
    \bottomrule
  \end{tabular}
\end{table*}

\begin{figure}[H]
  \centering
  \includegraphics[width=\linewidth]{f01_stance_ci.png}
  \caption{Stance distribution by platform with bootstrapped 95\% CIs.}
  \label{fig:stance}
\end{figure}

\subsection{Measures}
\label{sec:measures}
Three quantities recur and are defined explicitly.

\textbf{Stance consistency.} For a user $u$ with $n_u$ comments and $n_{u,s}$
comments of stance $s$, consistency is $\max_s n_{u,s}/n_u$ - the share of the user's
comments in their most frequent stance, where $1.0$ means the user never varies.
Computed over users with $\ge 3$ comments.

\textbf{Thread homophily.} With $T_u$ the set of distinct threads $u$ participates in
and $\mathrm{dom}(u)$ their dominant stance,
$H^{\text{thread}}_u = |\{t \in T_u : \text{majority}(t) = \mathrm{dom}(u)\}| / |T_u|$.
This measures co-presence, not interaction.

\textbf{Reply homophily and assortativity.} We build a directed user-user graph in
which an edge $u \rightarrow v$ exists whenever $u$ replies to a comment authored by
$v$, recovered from Reddit's \texttt{parent\_id} field (\texttt{t1\_} prefixes denote
replies to comments; self-replies are dropped). Reply homophily is the share of a
user's outgoing edges landing on same-stance users. Because that share depends on how
common each stance is among reply targets, we report it against the value expected
under random mixing and summarise the whole graph with Newman's assortativity
coefficient~\cite{newman2003}, where $0$ is random mixing and negative values indicate
cross-stance interaction.

\subsection{Analytical Methods}
Emotional tone is scored on every comment with VADER, TextBlob and
\texttt{twitter-roberta-base-sentiment}; RoBERTa, the instrument closest to the
human labels (Section~\ref{sec:sentval}), is primary, its continuous score being
$p(\text{positive}) - p(\text{negative})$. We report pairwise Cohen's $\kappa$ and
analyse the subset on which all three agree. Topics combine word frequency and
TF-IDF with LDA and NMF. Toxicity is scored on every comment with Detoxify
(\texttt{unbiased}), with Google's Perspective API on a balanced subsample of 600
comments as a cross-check. Stance classification uses a soft-voting ensemble (Logistic
Regression, calibrated Linear SVM, Random Forest) over TF-IDF features, evaluated
within and across platforms, with and without length matching, and explained by
\emph{permutation importance computed on the ensemble itself}. Because the corpus is
large, every significance test is paired with an effect size (Cram\'er's $V$,
rank-biserial $r$, or $\varepsilon^2$).

% ============================================================
\section{Results}
% ============================================================
\subsection{RQ1: Affective Tone and Toxicity}

\subsubsection*{Platform tone}
By RoBERTa, Reddit has a mean score of \textbf{$-0.536$} (71.1\% of comments
negative) and YouTube \textbf{$-0.169$} (45.0\%); a Mann-Whitney test confirms the
contrast ($p<0.001$; rank-biserial $r=+0.335$), and VADER agrees in direction
($-0.164$ against $+0.068$). This platform-level gap is not reproduced by human
labels (Section~\ref{sec:sentval}), which find the two platforms equally negative,
so we report it as an instrument result and draw no conclusion from it.

\textbf{Tone is bound to stance on Reddit, and only weakly on YouTube.} On Reddit
both partisan stances are strongly negative - Pro-Palestine ($-0.609$; 78.3\%
negative) and Pro-Israel ($-0.584$; 76.4\%) - while Neutral comments are far less so
($-0.265$; 43.1\%). On YouTube the three stances sit much closer together:
Pro-Palestine ($-0.220$), Pro-Israel ($-0.107$), Neutral ($-0.092$). The
stance-sentiment association is correspondingly far stronger on Reddit (Cram\'er's
$V=0.210$) than on YouTube ($V=0.073$), as Figure~\ref{fig:sentstance} shows, and
VADER gives the same contrast ($0.149$ against $0.039$). Negativity is thus the
signature of partisanship on Reddit, whereas on YouTube affect is only loosely tied
to political position.

\begin{figure*}[t]
  \centering
  \includegraphics[width=\textwidth]{f02_sentiment_by_stance.png}
  \caption{Mean RoBERTa sentiment by stance. Tone is politicised on Reddit and only
  weakly coupled to stance on YouTube.}
  \label{fig:sentstance}
\end{figure*}

%\subsubsection*{How much should we trust the sentiment tools?}
Agreement among the three methods is modest, and we report it rather than assume it
(Figure~\ref{fig:consensus}): across the full corpus, on Reddit $\kappa=0.291$
(VADER-RoBERTa), $0.217$ (VADER-TextBlob) and only $0.103$ (TextBlob-RoBERTa); on
YouTube the corresponding values are $0.413$, $0.291$ and $0.175$. TextBlob is
clearly the weakest instrument, and against human labels it performs close to
chance. Rather than average over disagreement, we isolate comments on which
\emph{all three} methods agree - 23.7\% of Reddit comments and 31.7\% of YouTube's.
On this subset the Reddit stance-sentiment association rises to $V=0.271$ while
YouTube stays low at $V=0.065$, so the substantive conclusion survives the
measurement concern. The three instruments also agree that Reddit is more negative
on this subset (57.2\% against 25.9\%), but agreement among instruments that may
share a blind spot is not validation, and the human labels do not show that gap.

\begin{figure*}[t]
  \centering
  \includegraphics[width=\textwidth]{f03_consensus.png}
  \caption{Left: pairwise inter-method agreement across the full corpus is modest,
  especially for TextBlob. Right: accuracy of each instrument against the human tone
  labels ($n=270$; dotted line marks chance for three classes).}
  \label{fig:consensus}
\end{figure*}

\subsubsection*{Toxicity by platform and stance}
Detoxify scores on all \textbf{1.38M} comments show Reddit higher on every
attribute (Table~\ref{tab:tox}): 17.8\% of Reddit comments exceed the conventional
$0.5$ toxicity threshold against 6.2\% on YouTube, and the largest platform effect
is identity attack ($r=-0.299$). Within platforms, both partisan camps score far
above Neutral on Reddit (mean toxicity $0.236$ and $0.229$ against $0.085$), with
the strongest stance-linked effects for identity attack ($\varepsilon^2=0.112$) and
threat ($\varepsilon^2=0.104$); on YouTube the ordering is the same but stance
effects are small ($\varepsilon^2 \le 0.018$). On the 600 comments also scored by
Perspective the two instruments rank comments similarly (Spearman $\rho=0.91$ on
Reddit and $0.69$ on YouTube for overall toxicity), and Perspective shows the same
platform ordering.

\begin{table}[H]
  \centering
  \small
  \caption{Mean Detoxify toxicity scores by platform (full corpus: 1{,}004{,}629
  Reddit and 378{,}267 YouTube comments).}
  \label{tab:tox}
  \begin{tabular}{lccc}
    \toprule
    Attribute & Reddit & YouTube & rank-biserial $r$ \\
    \midrule
    Toxicity        & \textbf{0.205} & 0.088 & $-0.266$ \\
    Severe toxicity & \textbf{0.003} & $<$0.001 & $-0.260$ \\
    Identity attack & \textbf{0.061} & 0.019 & $-0.299$ \\
    Insult          & \textbf{0.098} & 0.044 & $-0.243$ \\
    Threat          & \textbf{0.015} & 0.003 & $-0.267$ \\
    Obscenity       & \textbf{0.051} & 0.010 & $-0.202$ \\
    \bottomrule
  \end{tabular}
\end{table}

\subsubsection*{Are negativity and toxicity the same thing?}
That toxicity and negative sentiment are distinct constructs is well recognised -
toxicity classifiers target rudeness and disrespect rather than valence - but the
two are often used interchangeably in practice, and the size of the gap is rarely
quantified. In this corpus they are related but far from interchangeable
(Figure~\ref{fig:jointtox}). Sentiment and toxicity correlate moderately (Spearman
$\rho=-0.574$ on Reddit and $-0.484$ on YouTube with RoBERTa). Most toxic comments
are negative - 95.3\% on Reddit and 89.0\% on YouTube by RoBERTa, 75.3\% and 66.1\%
by VADER, and 74\% of the 19 toxic comments in the human-labelled validation sample
- but the converse fails badly: only \textbf{23.9\%} of negative Reddit comments and
\textbf{12.3\%} of negative YouTube comments are toxic (26.0\% and 13.1\% by VADER).
Toxicity is therefore a distinct subset of negative discourse, and treating
negativity as a proxy for toxicity would overcount it roughly fourfold.

\begin{figure*}[t]
  \centering
  \includegraphics[width=\textwidth]{f06_joint_toxicity.png}
  \caption{Negativity and toxicity overlap asymmetrically: most toxic comments are
  negative, but only about a quarter of negative Reddit comments and an eighth of
  negative YouTube comments are toxic (Detoxify, full corpus; negativity by RoBERTa
  and by VADER).}
  \label{fig:jointtox}
\end{figure*}

This distinction changes the stance story. Regressing toxicity jointly on sentiment
and stance, sentiment is a strong predictor (Reddit $b=-0.223$ per unit of RoBERTa
score), while the Pro-Israel versus Pro-Palestine difference is negligible on both
platforms ($b=-0.002$ on Reddit and $+0.002$ on YouTube, on a 0-1 scale). With a
million comments such differences are statistically detectable but substantively
nil, and they are equally small before sentiment is controlled ($-0.007$ and
$-0.009$). Attribute by attribute the camps differ only slightly, Pro-Israel
comments scoring marginally higher on identity attack and threat and Pro-Palestine
comments on insult and obscenity. What survives the control is the partisan-neutral
gap: comments of \emph{either} side are more toxic than neutral ones (Reddit
$b=-0.074$ $[-0.075, -0.073]$; YouTube $b=-0.026$), and VADER as the control gives
the same picture ($-0.120$ and $-0.033$).

\subsection{RQ2: Topics and Vocabulary}
The two platforms differ in vocabulary as well as tone. We report topic structure
descriptively: the labels below name the \emph{term families} a model groups
together and make no claim about who acted, who was responsible, or what any act
licenses. Recovering that would require statement-level narrative analysis, which we
do not perform.

\subsubsection*{A data-quality correction}
Our first pass produced a prominent Reddit ``community governance'' topic
(\textit{bot, rules, moderators, action, automatically}) and a Neutral class whose
most distinctive terms were almost entirely automated moderation text. Inspection
showed that automated accounts and the collapsed \texttt{[deleted]} placeholder
supply 5.93\% of the Reddit corpus, and that this contamination is strongly
asymmetric across stances: 14.2\% of Neutral and 6.5\% of Pro-Palestine comments,
but only 0.3\% of Pro-Israel comments, because subreddit rule text is itself
stance-inflected. Removing these accounts corpus-wide (not merely from user-level
analysis) both eliminates the spurious topic and shifts the stance composition
materially - Reddit's Pro-Israel share rises from 34.0\% to 36.1\%. All results in
this paper use the cleaned corpus.

\subsubsection*{Frequent and distinctive vocabulary}
After removing shared conflict terms, Reddit's most frequent words are
\textit{palestinians, jews, land, right, state, jewish, civilians, world, genocide},
and \textit{idf}. YouTube foregrounds \textit{god, allah, free, peace, bless, love,
children}, and \textit{support}. TF--IDF sharpens the stance contrast within each
platform: on Reddit, Pro-Palestine comments are distinguished by
\textit{genocide, land, right}, Pro-Israel comments by \textit{jews, jewish,
civilians}; on YouTube, Pro-Palestine by \textit{free, allah, land} and Pro-Israel
by \textit{god, bless, idf, stand}.

\subsubsection*{Topic structure}
LDA and NMF recover consistent structure on Reddit, which we label only by the
terms present:
\begin{itemize}[leftmargin=1.2em,itemsep=1pt,topsep=2pt]
  \item \textbf{territory and statehood terms:} israel, palestinians, gaza, state, west, bank, peace;
  \item \textbf{identity and land terms:} jews, jewish, land, arabs, arab, palestine;
  \item \textbf{conflict-actor and casualty terms:} hamas, civilians, idf, hostages, killed, children;
  \item \textbf{contested-framing terms:} genocide, propaganda, media, zionist, support;
  \item \textbf{conversational markers:} dont, think, know, youre, thats.
\end{itemize}
On YouTube the recovered families are markedly different:
\begin{itemize}[leftmargin=1.2em,itemsep=1pt,topsep=2pt]
  \item \textbf{Pro-Palestine slogan terms:} free, palestine, allah, river, sea, freedom;
  \item \textbf{Pro-Israel solidarity terms:} israel, stand, support, love, praying, india;
  \item \textbf{devotional terms:} god, bless, protect, jesus, amen, lord;
  \item \textbf{conflict-actor terms:} hamas, civilians, idf, hostages, blame, attack;
  \item \textbf{non-English tokens:} hai, ke, ko, ki, ya, indicating substantial
  Hindi/Urdu and Arabic-script participation.
\end{itemize}
Two observations follow. First, the same conflict is lexicalised through political
and legal vocabulary on Reddit and through devotional and slogan vocabulary on
YouTube. Second, YouTube's discourse is visibly multilingual and includes a large
South Asian participation stream that has no counterpart in our Reddit sample - a
composition difference that any cross-platform comparison of this conflict should
report, and which we return to when interpreting the stance-transfer results.

\subsubsection*{Topic-model quality}
We quantify topic quality with $c_v$ coherence rather than asserting it. NMF is
marginally more coherent than LDA on both platforms (Reddit $0.502$ vs $0.479$;
YouTube $0.431$ vs $0.419$), and Reddit topics are more coherent than YouTube's on
either model - consistent with YouTube's shorter, more formulaic comments providing
weaker co-occurrence signal. Because the two algorithms agree on the term families
reported above, the structure is not an artifact of a single modelling choice.

\subsubsection*{Comment length and readability}
Reddit comments average \textbf{49.3} words (median 27); YouTube comments average
\textbf{18.5} words (median 10). Flesch Reading Ease is correspondingly lower for
Reddit (\textbf{63.7}) than YouTube (\textbf{67.1}), i.e.\ Reddit comments are
slightly harder to read. We describe this strictly as \emph{readability}: Flesch is
a surface measure of syllable and sentence length and is not evidence of narrative
complexity or argumentative depth. Stance-level differences in readability are
negligible ($\varepsilon^2=0.008$ on Reddit, $0.012$ on YouTube). As noted in
Section~\ref{sec:data}, the length contrast is a property of the platforms' designs
rather than a finding about this conflict, and we use it below as a control rather
than as a result.


\subsection{RQ3: Polarization and Interaction Structure}

\subsubsection*{Engagement: a circular metric, and a non-circular replacement}
An OLS model predicting Reddit score from sentiment and stance has almost no
explanatory power ($R^2=0.009$). Controversial comments average a near-zero score
($0.17$) against $9.48$ for non-controversial ones. We stress that this is
\emph{definitional rather than behavioural}: Reddit sets the controversiality flag
precisely when a comment attracts substantial voting in both directions, and score is
upvotes minus downvotes, so a near-zero score follows arithmetically. We report it
only as a sanity check that the flag behaves as documented.

To ask whether controversy actually suppresses engagement we use a signal that is not
vote-derived: the number of direct replies a comment receives, recovered from the
reply graph. The answer reverses (Figure~\ref{fig:contro}). Controversial comments
attract \textbf{0.747} replies on average versus \textbf{0.335} for
non-controversial ones - more than twice as many ($p<0.001$, $r=-0.257$). Controversy
is penalised by the scoring mechanism while simultaneously generating more
conversation, a divergence invisible to vote-based measures alone.

\begin{figure}[H]
  \centering
  \includegraphics[width=\linewidth]{f08_controversy.png}
  \caption{Controversial comments score near zero by construction, yet receive more
  than twice as many replies.}
  \label{fig:contro}
\end{figure}

\subsubsection*{User consistency}
Among \textbf{44{,}593} active Reddit users ($\ge3$ comments) the dominant-stance
split is 23{,}978 Pro-Palestine, 15{,}670 Pro-Israel and 4{,}945 Neutral. Mean stance
consistency (Section~\ref{sec:measures}) is \textbf{0.803}: active users
overwhelmingly repeat the same stance. Individual-level consistency is therefore high
and is not in dispute.

\subsubsection*{Co-presence and interaction give opposite answers}
Measured by thread co-membership, the corpus looks like a textbook echo chamber: mean
thread homophily is \textbf{0.643} and 43.4\% of active users exceed $0.8$, with
Pro-Israel ($0.702$) and Pro-Palestine ($0.670$) far above Neutral ($0.329$).

The reply network tells a different story (Figure~\ref{fig:homcmp}). Across
\textbf{321{,}112} user-user reply edges, mean reply homophily is only
\textbf{0.406}. More importantly, that figure must be read against chance: because
52.0\% of reply targets are Pro-Palestine, a Pro-Palestine user replying at random
would already match 52\% of the time. Both partisan groups fall \emph{below} their
random-mixing baseline (Table~\ref{tab:mixing}), and the graph as a whole is
disassortative, with Newman $r=\mathbf{-0.184}$.

A marginal baseline alone would not settle this, because activity and popularity are
unevenly distributed across stances: if Pro-Palestine users are both more numerous
and more replied-to, some apparent cross-cutting would be mechanical. We therefore
repeat the test against a \emph{degree-preserving} null that permutes reply targets,
holding every user's out-degree and in-degree exactly fixed. Across 500 permutations
the null assortativity is centred on zero ($\bar{r}=+0.00004$, $\mathrm{sd}=0.0014$),
placing the observed $-0.184$ far outside the null distribution ($z=-128$). The
per-stance excesses survive the same test (Pro-Palestine $z=-126$, Pro-Israel
$z=-118$, Neutral $z=+51$). The cross-cutting pattern is therefore not an artifact
of the degree sequence. Pro-Israel users direct 64.4\% of
their replies at Pro-Palestine users; Pro-Palestine users direct 54.4\% of theirs at
Pro-Israel users (Figure~\ref{fig:mixing}). Among partisan reply pairs,
\textbf{62.4\%} cross the divide. Only Neutral users are mildly assortative
($+0.09$ above chance), and they are a small pool.

Crossing the divide changes how replies are written but not how toxic they are. At
the comment level, 280{,}097 Reddit replies connect two partisan comments. The
154{,}282 that answer the opposite stance are more often negative than same-stance
replies (85.5\% against 73.2\% by RoBERTa) and slightly more often toxic (20.5\%
against 19.1\%). Once the reply's own sentiment is controlled the toxicity
difference disappears (with VADER $b=+0.002$, with RoBERTa $b=-0.018$, on a 0-1
scale): cross-stance engagement is more negative, but not more toxic than its
negativity accounts for.

\begin{figure}[H]
  \centering
  \includegraphics[width=\linewidth]{f05_homophily_compare.png}
  \caption{Thread co-membership suggests clustering; the reply network does not.}
  \label{fig:homcmp}
\end{figure}

\begin{table}[H]
  \centering
  \small
  \caption{Reply homophily against the random-mixing baseline. Negative excess
  indicates cross-stance interaction.}
  \label{tab:mixing}
  \begin{tabular}{lccc}
    \toprule
    Replier stance & Observed & Expected & Excess \\
    \midrule
    Pro-Palestine & 0.406 & 0.520 & $-0.113$ \\
    Pro-Israel    & 0.307 & 0.425 & $-0.118$ \\
    Neutral       & 0.140 & 0.055 & $+0.085$ \\
    \midrule
    Overall       & 0.343 & 0.445 & $-0.102$ \\
    \bottomrule
  \end{tabular}
\end{table}

\begin{figure*}[t]
  \centering
  \includegraphics[width=\textwidth]{f04_reply_mixing.png}
  \caption{Who replies to whom. Partisan users reply across the divide more often
  than chance would predict; the stance reply graph is disassortative.}
  \label{fig:mixing}
\end{figure*}

The two measures are not contradictory; they measure different things. Users do
concentrate in threads dominated by their own side, but \emph{within} those threads
the replies they write disproportionately target the other side. Echo-chamber claims
resting on co-presence alone can therefore invert the conclusion drawn from
interaction. The cross-cutting direction itself agrees with De Francisci Morales
et al.~\cite{defrancisci2021}, who report the same preference - and a similar
asymmetry, with one camp more eager to reply across the divide - for US party
politics; recovering it in a conflict corpus indicates the pattern is not confined
to domestic partisanship.

\subsubsection*{Stance classification as a lexical diagnostic}
We treat stance classification not as a performance claim but as a diagnostic of how
differently the two platforms encode stance. A soft-voting ensemble reaches macro-F1
$0.640$ within Reddit and $0.470$ within YouTube (5-fold CV: $0.624\pm0.004$ and
$0.461\pm0.003$), and transfers poorly in both directions (Table~\ref{tab:ml}).

The obvious objection is that Reddit and YouTube comments differ sharply in length
(49.3 vs 18.5 words), so transfer failure might be a length artifact rather than a
vocabulary one. We tested this by matching the two length distributions bin-by-bin,
producing corpora with near-identical means (18.5 vs 18.4 words). The mean
within-minus-cross macro-F1 gap falls only from $0.140$ to $0.129$
(Figure~\ref{fig:transfer}). Length matching therefore does not explain the transfer
failure; the divergence is lexical.

\begin{table}[H]
  \centering
  \small
  \caption{Stance-classification macro-F1, raw and length-matched.}
  \label{tab:ml}
  \begin{tabular}{llcc}
    \toprule
    Train & Test & Raw & Length-matched \\
    \midrule
    Reddit  & Reddit  & 0.640 & 0.617 \\
    YouTube & YouTube & 0.470 & 0.475 \\
    Reddit  & YouTube & 0.403 & 0.420 \\
    YouTube & Reddit  & 0.427 & 0.414 \\
    \midrule
    \multicolumn{2}{l}{mean within$-$cross gap} & 0.140 & 0.129 \\
    \bottomrule
  \end{tabular}
\end{table}

\begin{figure}[H]
  \centering
  \includegraphics[width=\linewidth]{f07_transfer.png}
  \caption{Cross-platform transfer remains far below within-platform performance after length matching. Axis labels give train$\rightarrow$test, with \textbf{R} = Reddit and \textbf{Y} = YouTube: R$\rightarrow$R and Y$\rightarrow$Y are within-platform, R$\rightarrow$Y and Y$\rightarrow$R cross-platform.}%
  
  \label{fig:transfer}
\end{figure}

Permutation importance computed on the ensemble itself - so the explanation matches
the model actually evaluated - confirms the lexical account. Reddit's most
load-bearing terms are contested political nouns (\textit{hamas, israel, jews,
palestinians, zionists, terrorists, genocide}), while YouTube's mix slogan
vocabulary with channel-specific artifacts (\textit{free, palestina}, and
video-specific tokens such as \textit{podcast}). The presence of channel artifacts is
itself informative: part of what a YouTube stance classifier learns is which video a
comment sits under, which cannot transfer by construction.

% ============================================================
\section{Discussion}
% ============================================================
\subsection{Emotion Is Politicised on Reddit, Weakly Coupled on YouTube}
The platform-level tone gap depends on the instrument (Section~\ref{sec:sentval});
what is robust is \emph{where} emotion attaches to politics. On Reddit stance and
sentiment are appreciably associated ($V=0.210$, rising to $0.271$ among
unanimously-scored comments), with both partisan camps far more negative than
neutrals. On YouTube the association is weak ($V=0.073$). This is not predicted by
the observation that YouTube comments are short reactions: short reactions could
equally have been uniformly hostile. What the data show is that YouTube's affective
register - devotional, solidaristic or hostile - is spread far more evenly across
political positions, whereas on Reddit negativity is the marker of taking a side.

The size of that contrast should be read as an upper bound. Human validation
(Section~\ref{sec:validation}) shows our stance labels are considerably more accurate
on Reddit than on YouTube, and label noise attenuates exactly this kind of
association. Degrading Reddit's labels to YouTube's measured accuracy cuts its
association to $V=0.118$ - but leaves it 1.6 times YouTube's $0.073$, removing
off-topic comments outright leaves it 2.5 times YouTube's, and human labels for both
stance and tone put the ratio at about two. The qualitative claim, that affect is
politicised on one platform and only weakly coupled to stance on the other,
therefore stands; the quantitative gap is smaller than the raw figures suggest.
Bail et al.~\cite{bail2018} found that cross-cutting exposure
can intensify rather than soften partisan feeling, which is consistent with the
pattern we observe in Reddit's argumentative, high-exposure environment.

\subsection{Toxicity Is Not Simply Strong Negativity}
Treating negative sentiment as a proxy for toxic speech would be a substantial
error: only about one negative Reddit comment in four is toxic. The practical
implication for moderation research is that sentiment filters are poor triage tools
for toxicity. The finding also disciplines the partisan comparison. The two camps differ
slightly attribute by attribute, but on overall toxicity the difference between them
is negligible, with or without sentiment controlled, while the
partisan-versus-neutral gap remains. We therefore make no claim that either side is
the more toxic party; the defensible claim is that taking a side at all is associated
with more toxic language than remaining neutral. The same holds in direct exchange:
replies that cross the divide are more negative than same-stance replies, but not
more toxic once that negativity is accounted for.

\subsection{Co-Presence Is Not Interaction}
Our clearest methodological result is that the two standard operationalisations of
echo chambers disagree in sign. Thread co-membership yields homophily of $0.643$ and
supports an echo-chamber reading; the reply network yields $0.406$, below chance, and
a disassortative graph ($r=-0.184$) in which most partisan replies cross the divide.
Since homophily is defined on networks~\cite{newman2003}, the
reply-graph measure is the appropriate one, and it indicates that Reddit in this
corpus functions as a contested arena rather than a set of sealed chambers. That
conclusion is not itself new - De Francisci Morales et al.~\cite{defrancisci2021}
reached it for US party politics using a comparable degree-preserving null. What our
data add is the direct contrast, on identical users, with a co-presence measure that
points the other way, which suggests the persistent disagreement between
echo-chamber and no-echo findings is substantially a disagreement about what is
being measured. This does
not contradict high individual stance consistency ($0.803$): users are consistent in
\emph{what they argue} while still arguing largely with the other side. It aligns
with Shugars and Ha~\cite{shugars2025}, who found that Reddit's conversational
visibility supports more cross-cutting exchange than broadcast-oriented designs - our
reply graph is the network-level signature of exactly that affordance. Because
Cinelli et al.~\cite{cinelli2021} show echo-chamber strength varies by platform, we
caution that this result is specific to Reddit's threaded reply structure and should
not be generalised to feed-based platforms.

\subsection{Stance Language Is Platform-Specific}
That stance models transfer poorly across corpora is already established by Ng and
Carley~\cite{ng2022}, so we claim no novelty for the effect itself. What our design
adds is a control: the gap is essentially unchanged after matching comment-length
distributions, so it is not an artifact of Reddit's longer comments. The residual gap
is lexical - statehood, legal and atrocity vocabulary on Reddit against devotional
and slogan vocabulary on YouTube - and is partly driven by video-specific tokens that
cannot transfer even in principle. A further compositional difference matters here:
YouTube's comment stream is visibly multilingual, with substantial Hindi/Urdu and
Arabic-script participation absent from our Reddit sample, so the two corpora differ
in speaker population as well as vocabulary. The implication is for study design
rather than modelling: conclusions about stance drawn from a single platform should
not be assumed to describe the discourse elsewhere.

\subsection{Limitations}
Several limitations qualify these results. First, stance labels were assigned by a
large language model whose accuracy we measure rather than assume
(Section~\ref{sec:validation}): $0.836$ on Reddit but $0.618$ on YouTube against a
three-annotator gold standard, with roughly one comment in six carrying a stance label
that humans judge irrelevant. Cross-platform comparisons therefore carry unequal
attenuation; our simulation shows the headline contrast survives it at about half its
raw size, but YouTube stance-conditional effects should be read as directional.
We release a relevance classifier that removes most of this content, and show the reported comparisons survive it, but do not adopt it as the primary corpus because it also discards genuine Neutral comments. The clearest routes to tightening these estimates are specific rather than general, and we measure both: supplying video-level context to the YouTube annotator raises stance accuracy by roughly eight points but degrades relevance judgement, so it must be paired with an explicit relevance decision ahead of
stance assignment. Second, toxicity is measured by automated classifiers with no
human toxicity labels behind them. Detoxify agrees with Perspective on the 600
comments both scored ($\rho=0.91$ on Reddit, $0.69$ on YouTube), but both are
trained on English-language rater data, and the sentiment validation shows that
automated instruments under-read negativity in YouTube's informal comments; the
platform toxicity gap may be overstated for the same reason, although its direction
holds under both classifiers. Third, YouTube lacks usable comment timestamps, so
its temporal anchor is the video publish date and no within-video temporal dynamics
can be recovered. Fourth, the reply-network analysis is Reddit-only: YouTube's flat
comment structure provides no comparable reply graph, so the central RQ3 result is
single-platform. Fifth, coverage is partial - five subreddits and a fixed YouTube
query set - and the two platforms differ in language composition, so neither corpus is
a census nor are their speaker populations matched. Finally, sentiment instruments
agree only modestly with each other ($\kappa$ as low as $0.103$), and against human
labels the best of them reaches an accuracy of $0.68$. We therefore rest RQ1 on the
stance-tone association, which human labels reproduce, and not on platform-level
tone, which they do not.

% ============================================================
\section{Conclusion}
% ============================================================
Across 1.38M stance-labelled comments in a matched window, the same conflict is
organised differently by two platforms - but the differences that matter are not the
ones visible from platform design. Emotional tone is politicised on Reddit and only
weakly coupled to stance on YouTube, a contrast that human labels reproduce.
Toxicity is a distinct and much rarer phenomenon than negativity, and the two
partisan camps are equally toxic while both exceed neutral comments. Most consequentially, Reddit's users cluster by thread but
argue across stance lines: the reply network is disassortative, and 62.4\% of
partisan replies cross the divide, so echo-chamber conclusions drawn from co-presence
invert when interaction is measured directly. Stance remains encoded in
platform-specific vocabulary, and we show this is not an artifact of comment length.

We also report what the labels will and will not support. Three blind annotators put
stance accuracy at $0.836$ on Reddit and $0.618$ on YouTube, and show that about one
comment in six carries a stance label humans would call irrelevant. Degrading Reddit's
labels to YouTube's measured quality shrinks the emotional-tone contrast without
eliminating it, and human tone labels reproduce that contrast but not a platform
difference in overall negativity. The platform differences reported here are therefore real but smaller
than the raw figures imply, and the clearest next step is not a further analysis of
this corpus but a better-annotated one.

% ============================================================
\section*{Data and Code Availability}
% ============================================================
The labelling pipeline, analysis code, revision scripts, the three completed human
validation annotation sets, the sentiment labels of both annotators and the
adjudicator, the scoring scripts and the video-level metadata are
available at \url{https://github.com/akshat3144/israel-hamas-discourse-analysis}.

\bibliographystyle{ieeetr}
\bibliography{refs}

\end{document}
